\clearpage
\renewcommand{\sectioncolor}{nano}
\section{What AI actually adds here --- and one thing to be careful of}
\markboth{AI, honestly}{}

\noindent\textbf{\color{nano}The biology question:} \emph{everyone says AI will transform genomics.
For our pipeline, specifically, what does it do?}

\subsection{The honest ledger}

\begin{center}\small
\begin{tabular}{@{}>{\raggedright\arraybackslash}p{3.1cm}>{\raggedright\arraybackslash}p{5.5cm}c>{\raggedright\arraybackslash}p{6.4cm}@{}}
\toprule
\rowcolor{nano!13}
\textbf{\color{nano}Step} & \textbf{\color{nano}What exists} & &\textbf{\color{nano}Use to us}\\
\midrule
Base calling & DeepConsensus (Baid et al., \textit{Nat Biotechnol} 2023) --- transformer on raw
 PacBio signal; Dorado/Bonito for nanopore & \pt &
 Real but not ours: these are long-read tools. Illumina base-calling gains are now small --- our own
 handbook §T7.4 traces Bustard $\to$ Ibis $\to$ AYB, 4.0\,\% $\to$ 2.8\,\% error.\\
Variant calling & DeepVariant (Poplin et al., \textit{Nat Biotechnol} 2018) & \pt &
 Mature and excellent --- but relevant only if we move from \emph{assembly} to \emph{resequencing
 against a reference}.\\
\rowcolor{bio!8}
Gene identification & \textbf{Protein language models (ESM-2; Lin et al., \textit{Science} 2023)} & \yes &
 \textbf{Directly useful.} Embeddings separate orthologues from paralogues far better than percent
 identity --- this is the \S7 problem.\\
Genome-scale priors & Nucleotide Transformer (\textit{Nat Methods} 2024); Evo (\textit{Science} 2024) & \pt &
 Promising, and see the warning below.\\
Structure & AlphaFold2 (Jumper et al., \textit{Nature} 2021) & \pt &
 Useful for interpreting a novel QRDR substitution; not part of routine screening.\\
\rowcolor{bio!8}
Analysis assistance & what we already do --- AI as shell and scripting assistant & \yes &
 \textbf{Already working, and our \texttt{METHODS} §7 states the right rule: every number traces to a
 deterministic re-runnable command.}\\
\midrule
\multicolumn{4}{@{}l}{\textbf{\color{alert}What AI cannot do for us}}\\
\multicolumn{4}{@{}>{\raggedright\arraybackslash}p{16.6cm}@{}}{%
 \no~reduce photon shot noise \quad
 \no~span a 2{,}600 bp repeat with a 233 bp read \quad
 \no~detect a failed tile-cycle better than the 50-line script in \S10}\\
\bottomrule
\end{tabular}
\end{center}

\subsection{The one genuinely deep idea}

Assembly and error correction are both a search for the most probable genome given the reads:
\[
\hat{X} \;=\; \arg\max_{X}\;\underbrace{P(\text{reads}\mid X)}_{\text{what the machine measured}}\;\times\;
\underbrace{P(X)}_{\text{what we expected beforehand}}
\]
Coverage supplies the first term. \textbf{A good prior supplies the second --- and a sharp enough
prior substitutes for coverage.} That is why reference-guided analysis needs far less data than
\emph{de novo} assembly: the reference \emph{is} a prior.

\begin{plain}
If you already know roughly what the answer should look like, you need less evidence to confirm it.
That is true in the lab and it is true in mathematics. A genome language model is an attempt to
supply that ``roughly know'' for organisms where we have no close reference.
\end{plain}

\begin{alertbox}[title={The warning that matters specifically for AMR work}]
A learned prior pulls the estimate \textbf{toward what was common in its training data}.
\begin{center}
\begin{tikzpicture}[font=\footnotesize,
  b/.style={rounded corners=4pt,draw=#1,line width=1.1pt,fill=#1!8,inner sep=8pt,
            text width=7.3cm,align=left,minimum height=2.9cm}]
 \node[b=bio] (a) at (-4.25,0)
   {\textbf{\color{bio}Where the prior helps}\\[4pt]
    The isolate belongs to a common, well-sampled lineage. Most of the genome is exactly what the model has seen thousands of
    times. The prior is right, and it saves data.};
 \node[b=alert] (b) at (4.25,0)
   {\textbf{\color{alert}Where the prior hurts}\\[4pt]
    A \emph{novel} resistance allele. A new SCC\emph{mec} junction. A rearranged plasmid.\\[3pt]
    {\scriptsize The model has never seen it, so it will confidently smooth it into the nearest thing
    it has seen.}};
\end{tikzpicture}
\end{center}
\vspace{-2pt}
\begin{center}\bfseries\itshape
AMR research is the study of what is not yet common.\\
The novel findings are precisely the ones a learned prior will most reliably damage.
\end{center}
\vspace{3pt}
Our \texttt{METHODS} §7 already says ``no result is a model assertion''. The extension worth adopting
as team policy: \textbf{a learned prior is a model assertion smuggled into the data layer.} If we ever
adopt AI-assisted assembly or error correction, any novel call must be re-verified against the raw
reads without the model in the loop.
\end{alertbox}
